% Compile with xelatex
\documentclass[10pt,a4paper]{article}

\usepackage{fontspec}
\setmainfont{Malgun Gothic}[Ligatures=TeX]
\setsansfont{Malgun Gothic}[Ligatures=TeX]

\usepackage[top=2.0cm, bottom=2.0cm, left=2.5cm, right=2.5cm]{geometry}
\usepackage{microtype}
\usepackage{amsmath}
\usepackage{fancyhdr}

% ─── Conference running header (as in template) ─────────────
\pagestyle{fancy}
\fancyhf{}
\renewcommand{\headrulewidth}{0.4pt}
\fancyhead[L]{\small 2026년 한국생산운영관리학회 추계학술대회 발표 초록집}
\fancyfoot[C]{\small -- \thepage{} --}

\setlength{\parindent}{0pt}
\setlength{\parskip}{3pt}
\setlength{\emergencystretch}{1em}

% ============================================================
\begin{document}

% ============================================================
% PAGE 1 — Korean
% Primary title: Korean; secondary (small italic): English
% ============================================================
\begin{center}
{\Large\textbf{국내 관광명소 온라인 리뷰에서 방문 맥락과\\[2pt]
리뷰 결과에 관한 실증분석}}

\vspace{12pt}

{\normalsize *소리아노 에르난데스 호르헤 아단 (호서대학교 데이터사이언스학과,
대학원생, adan\_25@naver.com)}

\end{center}

\vspace{10pt}

\noindent\textbf{\large 초록}

\vspace{4pt}

\noindent
본 연구는 한국 관광지에 대한 네이버 플레이스 온라인 리뷰에서, 플랫폼이 구조화하여
수집한 방문 맥락(동행 유형, 방문 목적, 대기 여부, 예약 여부)이 리뷰 결과에 어떻게
반영되는지를 분석하고자 한다. 선행 연구는 리뷰 텍스트의 감성이나 주제 구조에
집중해 왔으며, 플랫폼이 제공하는 방문 맥락 메타데이터를 공동 예측 집합으로 활용하여
리뷰 결과를 설명하려는 시도는 상대적으로 부족하였다. 또한 많은 연구가 단일 관광지나
단일 범주에 국한되어, 다양한 유형의 관광지를 포괄하는 리뷰 수준의 비교 분석이
부족한 실정이다. 이에 본 연구는 두 가지 연구 질문을 다룬다. 첫째, 방문 맥락과
작성자 특성이 리뷰어의 선언된 만족도(별점)에 어떻게 영향을 미치는가. 둘째, 방문
맥락과 작성자 특성이 재방문 여부에 어떻게 반영되는가. 두 결과 변수는 관측 불가능한
만족도의 서로 다른 흔적으로 간주하며, 별점은 선언된 지표, 재방문 표식은 과거
경험에 기반한 행동 지표로 해석한다. 분석 대상은 방문객 규모를 기준으로 한 층화 추출로 선정된 서울·부산 등 29개 주요
관광지의 한국어 리뷰 14{,}401건이며, 별점에 대한 순서형 로지스틱 회귀와 재방문
여부에 대한 이항 로지스틱 회귀를 활용한다. 두 모형은 방문 맥락, 작성자 특성, 리뷰
길이, 장소 통제를 공통 예측 변수로 사용한다. 본 연구는 플랫폼이 구조화한 방문
맥락을 리뷰 수준의 공동 예측 집합으로 활용하여, 선언된 만족도와 행동적 만족도의
흔적을 동시에 설명하는 분석 틀을 제시한다는 데 의의가 있다.

\vspace{6pt}

\noindent\textbf{주제어 (Keywords):} 리뷰, 방문 맥락, 재방문, 만족도,
로지스틱 회귀

% ============================================================
% PAGE 2 — English
% Primary title: English; secondary (small italic): Korean
% ============================================================
\newpage

\begin{center}
{\Large\textbf{An Empirical Analysis of Visit Context and\\[2pt]
Review Outcomes in Online Reviews of Korean Tourist Attractions}}

\vspace{12pt}

{\normalsize *Jorge Adan Soriano Hernandez (Department of Data Science,
Hoseo University, Graduate Student, adan\_25@naver.com)}

\end{center}

\vspace{10pt}

\noindent\textbf{\large Abstract}

\vspace{4pt}

\noindent
This study aims to examine how platform-structured visit context (companion
type, visit purpose, whether the reviewer waited in line, and whether a
reservation was made) is reflected in the outcomes of online reviews of
Korean tourist attractions on Naver Place. Prior work has focused on review
text, sentiment, or topic structure, and the use of platform-provided
visit-context metadata as a joint predictor set to explain review-level
outcomes remains limited. Moreover, most existing studies are restricted to
a single attraction or category, leaving a gap in review-level comparative
analysis across heterogeneous destinations. This study addresses two
research questions. First, how do visit context and reviewer characteristics
shape the reviewer's declared satisfaction, expressed as the star rating?
Second, how do visit context and reviewer characteristics relate to the
revisit indicator? Both outcomes are interpreted as different observable
traces of an unobservable satisfaction construct, with the star rating as a
declared signal and the revisit flag as a behavioural one grounded in
prior experience. The dataset consists of 14{,}401 Korean-language reviews
from 29 major Korean tourist attractions in and around Seoul and Busan,
selected via stratified sampling based on visitor volume. An ordered
logistic regression of the star rating and a binary logistic regression of
the revisit indicator are employed, both using visit context, reviewer
characteristics, review length, and place controls as common predictors.
This study contributes by proposing an analytic framework that uses
platform-structured visit context as a joint predictor set at the review
level, explaining both declared and behavioural traces of satisfaction
simultaneously.

\vspace{6pt}

\noindent\textbf{주제어 (Keywords):} review, visit context, revisit,
satisfaction, logistic regression

\end{document}
